% -*- mode: latex; coding: utf-8 -*-
% pxTFG -- Standalone user manual and public API reference
% P3CTeX / Scvria. Compile with pdflatex doc/pxTFG.tex from tex/.
\documentclass[11pt,a4paper]{article}
\usepackage[T1]{fontenc}
\usepackage[utf8]{inputenc}
\usepackage{lmodern}
\usepackage[english]{babel}
\usepackage{pxTFG}
\usepackage[colorlinks=true,linkcolor=blue,urlcolor=blue]{hyperref}
\title{The \texttt{pxTFG} package\\\large UOC final-work metadata and front matter}
\author{Scvria\\P3CTeX / UOC PEC Repository}
\date{2026/10/10 -- Version 0.1}
\begin{document}
\maketitle
\begin{abstract}
\texttt{pxTFG} adds UOC final-project semantics to P3CTeX without
changing its pxGDX visual style. It is independent from the generic
\texttt{pxMEM} report and content-reuse package.
\end{abstract}
\tableofcontents
\setlength{\emergencystretch}{4em}

\section{Enable TFG}
The TFG module loads via pxCORE when the \texttt{TFG} class option is set:
\begin{verbatim}
\documentclass[TFG,PR,CAT]{P3CTeX}   % standalone PR
\documentclass[TFG,MEM,CAT]{P3CTeX}  % final memory
\end{verbatim}
Neither \texttt{TFG} nor \texttt{MEM} implies the other. Select them at
class-loading time; \verb|\PECTeXconfig{TFG}| and
\verb|\PECTeXconfig{MEM}| do not load modules retroactively.
The existing PR, PAC, PEC and CA profiles retain their semantics.

\section{Metadata}
\texttt{pxTFG} reuses P3CTeX data rather than creating duplicates.
Use the normal class data keys \texttt{student} (author),
\texttt{grade} (degree), \texttt{subj-fullname} (project title),
\texttt{ca-fullname}/\texttt{ca-shortname} (deliverable), and
\texttt{date} (document date). Use the class language option
\texttt{CAT}, \texttt{ES} or \texttt{EN}; there is no TFG language field.

\subsection{Configuration: \texttt{\textbackslash pxTFGsetup}}
\verb|\pxTFGsetup{<key=value,...>}| stores four TFG-only fields:
\begin{description}
\item[\texttt{tutor}] Final-work supervisor.
\item[\texttt{area}] Work area/specialisation.
\item[\texttt{credits}] Credit count (ECTS).
\item[\texttt{licence}] Licence text for the work card.
\end{description}
\begin{verbatim}
\PECTeXconfig{data={
  student={Estudiant},grade={Enginyeria Informatica},
  subj-fullname={Projecte},ca-fullname={Proposta},
  ca-shortname={PR1},date={\today}
}}
\pxTFGsetup{
  tutor={Tutora},area={Computacio},
  credits={12},licence={CC BY}
}
\end{verbatim}

\section{Public commands}
\subsection{\texttt{\textbackslash pxTFGworkcard}}
\verb|\pxTFGworkcard| prints the UOC final-work information card,
using both the shared class data and the four TFG fields. Its labels
follow the document language: Fitxa del Treball Final, Ficha del
Trabajo Final or Final Project Information. It creates an unnumbered
section in TFG PR documents or an unnumbered chapter in TFG+MEM
documents, with a table-of-contents entry.

\subsection{\texttt{\textbackslash pxTFGabstract}}
\verb+\pxTFGabstract{<language>}{<text>}+ registers one abstract for
a particular language. Re-registering the same language replaces its text.
It does not print the abstract and does not change Babel's document language.
\begin{verbatim}
\pxTFGabstract{CAT}{Resum del treball.}
\pxTFGabstract{ES}{Resumen del trabajo.}
\pxTFGabstract{EN}{Project abstract.}
\end{verbatim}

\subsection{\texttt{\textbackslash pxTFGabstracts}}
\verb|\pxTFGabstracts| prints the abstracts already registered in
CAT, ES, EN order, skipping absent entries. Headings are Resum,
Resumen and Abstract; each appears in the table of contents.

\subsection{\texttt{\textbackslash pxTFGfrontmatter}}
\verb|\pxTFGfrontmatter| prints the work card followed by the
registered abstracts. It does not print a cover or TOC; those belong
to P3CTeX. Place it in the body near the beginning of the document.

\section{Standalone TFG PR example}
\begin{verbatim}
\documentclass[TFG,PR,CAT]{P3CTeX}
\PECTeXconfig{data={
  student={Estudiant},grade={Enginyeria Informatica},
  subj-fullname={Projecte},ca-fullname={PR1},
  ca-shortname={PR1}
}}
\pxTFGsetup{tutor={Tutora},area={Computacio},
  credits={12},licence={CC BY}}
\begin{document}
\pxTFGworkcard
\section{Objectius}
Text de la PR1.
\end{document}
\end{verbatim}
A normal PR remains independently compilable. Potentially reusable
sections can be surrounded by selection comments as described in the
\texttt{pxMEM} manual; no TFG-to-MEM dependency is introduced.

\section{Final TFG memory example}
\begin{verbatim}
\documentclass[TFG,MEM,CAT]{P3CTeX}
\PECTeXconfig{cover=true,toc=false,data={
  student={Estudiant},grade={Enginyeria Informatica},
  subj-fullname={Projecte},ca-fullname={Memoria},
  ca-shortname={MEM}
}}
\pxTFGsetup{tutor={Tutora},area={Computacio},
  credits={12},licence={CC BY}}
\pxTFGabstract{CAT}{Resum.}
\pxTFGabstract{ES}{Resumen.}
\pxTFGabstract{EN}{Abstract.}
\begin{document}
\pxTFGfrontmatter
\tableofcontents
\pxMEMsource{PR1}{../PR1/PR1.tex}
\chapter{Introduccio}
\pxMEMinclude{PR1}{objectius}
\pxMEMannex{Evidencies}{Material addicional.}
\end{document}
\end{verbatim}
\section{Public API summary}
\begin{description}
\item[\texttt{\textbackslash pxTFGsetup\{keys\}}] Set TFG-only metadata.
\item[\texttt{\textbackslash pxTFGabstract\{lang\}\{text\}}] Store an abstract.
\item[\texttt{\textbackslash pxTFGworkcard}] Print the UOC work card.
\item[\texttt{\textbackslash pxTFGabstracts}] Print registered abstracts.
\item[\texttt{\textbackslash pxTFGfrontmatter}] Print card plus abstracts.
\end{description}
\end{document}
